\documentclass[11pt]{amsart}
\usepackage[T1]{fontenc}
\usepackage[margin=1.02in]{geometry}
\DeclareFontShape{OT1}{cmr}{bx}{n}{<-> cmbx10}{}
\DeclareMathSizes{8}{8}{8}{8}
\DeclareMathSizes{9}{9}{8}{8}
\DeclareMathSizes{10.95}{10.95}{8}{8}
\newtheorem{theorem}{Theorem}
\title{Boolean boundedness need not transfer to tropical Datalog}
\date{November 27, 2026}
\begin{document}
\begin{abstract}
A fixed linear monadic Datalog program stabilizes on every
Boolean input by round two, yet needs arbitrarily many naive
rounds over the nonnegative integer tropical semiring. Its
short Boolean proof uses a guarded input fact twice on separate
branches. This duplication can make it more expensive than a
longer tropical proof. This answers the boundedness-transfer
question of Fan, Koutris and Roy. A separate target-only example
has a constant-depth Boolean circuit and worst-case tropical
depth of order the square of a logarithm.
\end{abstract}
\maketitle

\section{A three-rule program}
Fan, Koutris and Roy \cite[Sec.~4, p.~93:10]{fkr} ask whether a
program can be Boolean-bounded and unbounded over the tropical
semiring. Boundedness means that naive synchronous evaluation of
a designated intensional predicate reaches its final values in
a fixed number of rounds independent of the finite database.
We use the nonnegative integer tropical semiring
\[
      (\mathbb N\cup\{\infty\},\ \min,\ +,\ \infty,\ 0).
\]
A missing extensional fact has weight \(\infty\). Numeric zero
is the weight of a present fact of multiplicative identity.
Intensional facts start at the additive identity in round zero.
Each round evaluates all rules using the preceding intensional
values. Boolean evaluation similarly uses disjunction and
conjunction with all intensional facts false in round zero.

\begin{theorem}
There is a fixed linear monadic program with acyclic rule bodies
for which all intensional relations stabilize over Boolean inputs
by round two. Its target requires arbitrarily many naive rounds
over the nonnegative integer tropical semiring.
\end{theorem}

\begin{proof}
Use unary extensional predicate \(G\), binary extensional
predicate \(E\), and unary intensional predicates \(U,T\).
The target is \(T\). The rules are
\[
\begin{aligned}
 U(x)&\ \leftarrow\ G(x),\\
 T(x)&\ \leftarrow\ G(x),U(x),\\
 T(x)&\ \leftarrow\ G(x),E(x,y),T(y).
\end{aligned}
\]
Each body has at most one intensional atom and is acyclic as a
hypergraph. In the Boolean semiring \(U=G\) after round one.
Every derivation of \(T(x)\) requires \(G(x)\), and the second
rule derives exactly those facts by round two. Both relations
have reached their final values, on every finite Boolean input.

Over tropical weights take the directed path \(0\to1\to
\cdots\to n\), where each edge is present with numeric weight
zero. Let \(G(i)\) have weight \(w_i=3^{n-i}\). All finite
derivations of \(T(i)\) follow the path to a stopping vertex
\(j\ge i\), then use the second rule and expand \(U(j)\) by the
first. Each such proof tree has height \(j-i+2\), with input
leaves at height zero, and cost
\[
       C(i,j)=\sum_{k=i}^{j-1}w_k+2w_j .
\]
The two copies of \(G(j)\) occur on separate branches, one under
\(U(j)\). No label repeats on a root-to-leaf path. Thus these
are also tight proof trees in the source sense
\cite[Sec.~2.4]{fkr}. By induction on round number, the round
\(r\) value is the minimum cost among proof trees of height
at most \(r\), with minimum of an empty set \(\infty\).
The rule recurrence gives exactly the same statement.

When \(j<n\), the drop in cost on extending the stopping vertex
by one is \(w_j-2w_{j+1}=w_{j+1}>0\). Thus \(T(0)\) first
attains its final value when the path reaches \(n\), in round
\(n+2\). This goes to infinity with \(n\).
\end{proof}

For example, four edges give weights \(81,27,9,3,1\). At
\(T(0)\) the successive finite values in rounds two through
six are \(162,135,126,123,122\). All corresponding Boolean
facts already stabilized by round two.

Source Corollary~4.7 transfers Boolean boundedness to absorptive
semirings with idempotent multiplication. Numeric addition in
the tropical semiring is not idempotent. The second rule
charges the two occurrences that Boolean conjunction absorbs.
The theorem concerns naive rounds of a fixed program. Existing
size-dependent convergence bounds for grounded systems are
compatible \cite{frisk}.

\section{Companion separation for circuit depth}
A separate construction isolates circuit depth of a single
target. Use binary extensional predicates \(F,E\) and binary
intensional predicates \(W,R,P\), where \(P\) is the target:
\[
\begin{aligned}
 W(x,y)&\ \leftarrow\ F(x,y),\\
 R(x,y)&\ \leftarrow\ E(x,y),\\
 R(x,y)&\ \leftarrow\ E(x,z),R(z,y),\\
 P(x,y)&\ \leftarrow\ F(x,y),W(x,y),\\
 P(x,y)&\ \leftarrow\ F(x,y),R(x,y).
\end{aligned}
\]
Boolean \(P\) equals \(F\) from round two, and its provenance
circuit for each target fact is a single input wire or a
constant. Internal \(R\) need not be Boolean-bounded.

In the tropical case \(R\) gives positive-length shortest-path
cost, and the target is
\[
       P(x,y)=\min\{2F(x,y),\,F(x,y)+R(x,y)\}.
\]
Fix a complete directed structure on \(N\) vertices and a single
\(F(s,t)\) fact, \(s\ne t\). Set its weight to one. Give each
\(E\) fact weight zero for a selected edge and two otherwise,
so the structure is fixed for all of these annotations.
Then \(P(s,t)\) is one exactly when there is a selected directed
path, and is two otherwise.

Suppose a fan-in-two min-plus circuit computes \(P(s,t)\) for
all annotations of this structure, with constants restricted
to nonnegative integers or \(\infty\). For each gate output
\(z\), keep the Boolean tests \(B_0(z)=[z\le0]\) and
\(B_1(z)=[z\le1]\). At a minimum gate each test is an OR. At
an addition gate with children \(a,b\), the tests are
\[
 B_0(a)\land B_0(b),\qquad
 \bigl(B_0(a)\land B_1(b)\bigr)
 \lor\bigl(B_1(a)\land B_0(b)\bigr).
\]
These identities hold on nonnegative integers with infinity.
At leaves the tests are constants or selected-edge inputs.
This turns tropical circuit depth \(d\) into monotone Boolean
depth at most \(2d\). Identify the variables for each pair of
opposite arcs to obtain undirected \(s\)-\(t\) connectivity.
The Karchmer--Wigderson lower bound \cite{kw} gives depth
\(\Omega(\log^2 N)=\Omega(\log^2 m)\) for the tropical circuit,
where \(m\) is the number of input facts.

For the upper bound, nonnegative weights permit removal of
cycles from a shortest path between distinct vertices.
For a positive-length shortest walk from a vertex to itself,
one can retain a shortest simple cycle. In both cases, if a
walk exists, some optimum uses at most \(N\) edges. Let \(A\)
be the min-plus adjacency matrix and \(I\) the min-plus
identity. Compute \(A(I\oplus A)^k\), with \(k\) a power of
two at least \(N-1\), using repeated squaring. The first factor
forces a real edge, preserving positive length even on the
diagonal. Each product costs depth \(O(\log N)\) at fan-in
two, and the number of products is \(O(\log N)\). Forming
\(P\) from \(R\) adds constant depth. The active domain has
\(O(m)\) vertices. With the standard path-doubling method
\cite[Sec.~5]{fkr}, this proves the matching \(O(\log^2 m)\)
upper bound. The worst-case depth claim is for this target
and this circuit model. The general depth classification of
Datalog remains a separate question.

\section*{Verification and technical disclosure}
The Python replay checks exact path costs for 200 lengths,
all Boolean databases on one to three vertices, and finite
controls for the circuit conversion and matrix formula.
A partial Lean~4 certificate checks the worked path trace,
the strict adjacent-cost comparison, two finite natural
threshold identities, and a Boolean absorption identity.
The general semantics and circuit lower bound are established
in the written proof.

Internally reviewed, not peer reviewed. Under human direction,
ChatGPT Astra~6.2 (OpenAI) generated the discovery proof and
research packet. OpenAI Codex independently checked the
mathematics. GPT-6.2 Sol (OpenAI) checked the literature and
prepared the manuscript, additional exact replay, partial Lean
certificate and illustrations. DannyExperiments maintains the
repository and is the preferred citation identity. No AI system
is proposed as an author.

\begin{thebibliography}{9}
\raggedright
\bibitem{fkr} A. Z. Fan, P. Koutris and S. Roy,
\emph{Circuits and formulas for Datalog over semirings},
Proc. ACM Manag. Data 3(2), PODS (2025), Art. 93.
doi:10.1145/3725230.
\bibitem{frisk} S. Frisk, S. Im, P. Koutris, B. Moseley,
H. Q. Ngo, K. Pruhs and H. Zhao,
\emph{Optimal convergence of iterative methods for Datalog},
arXiv:2610.04062v1, 2026.
\bibitem{kw} M. Karchmer and A. Wigderson,
\emph{Monotone circuits for connectivity require super-logarithmic depth},
SIAM J. Discrete Math. 3(2) (1990), 255--265.
\end{thebibliography}
\end{document}
